\documentclass{article}

% Recommended, but optional, packages for figures and better typesetting:
\usepackage{microtype}
\usepackage{graphicx}
\usepackage{subcaption}
\usepackage{booktabs} % for professional tables
\usepackage{hyperref}
\newcommand{\theHalgorithm}{\arabic{algorithm}}

% Use the accepted style for professional author layout
\usepackage[accepted]{icml2026}

% Global spacing adjustments to optimize layout and fit strict 8-page main paper constraints
\setlength{\textfloatsep}{5pt plus 1pt minus 1pt}
\setlength{\intextsep}{5pt plus 1pt minus 1pt}
\setlength{\floatsep}{5pt plus 1pt minus 1pt}
\setlength{\dbltextfloatsep}{5pt plus 1pt minus 1pt}
\setlength{\dblfloatsep}{5pt plus 1pt minus 1pt}
\setlength{\abovedisplayskip}{3pt plus 1pt minus 1pt}
\setlength{\belowdisplayskip}{3pt plus 1pt minus 1pt}
\setlength{\abovedisplayshortskip}{3pt plus 1pt minus 1pt}
\setlength{\belowdisplayshortskip}{3pt plus 1pt minus 1pt}
\renewcommand{\baselinestretch}{0.98} % Tiny baseline stretch compression (less than 2% for strict formatting adherence)

\usepackage{amsmath}
\usepackage{amssymb}
\usepackage{mathtools}
\usepackage{amsthm}

\usepackage[capitalize,noabbrev]{cleveref}

%%%%%%%%%%%%%%%%%%%%%%%%%%%%%%%%
% THEOREMS
%%%%%%%%%%%%%%%%%%%%%%%%%%%%%%%%
\theoremstyle{plain}
\newtheorem{theorem}{Theorem}[section]
\newtheorem{proposition}[theorem]{Proposition}
\newtheorem{lemma}[theorem]{Lemma}
\newtheorem{corollary}[theorem]{Corollary}
\theoremstyle{definition}
\newtheorem{definition}[theorem]{Definition}
\newtheorem{assumption}[theorem]{Assumption}
\theoremstyle{remark}
\newtheorem{remark}[theorem]{Remark}

\icmltitlerunning{\smash{Preserving ReLU Sparsity via Layer-wise Calibration}}

\begin{document}

\twocolumn[
\icmltitle{Less is More: Preserving ReLU Sparsity via Layer-wise \\
  Scaling-only Activation Calibration in Multi-Task Model Merging}

\icmlsetsymbol{equal}{*}

\begin{icmlauthorlist}
\icmlauthor{The Minimalist Agent}{equal,inst}
\end{icmlauthorlist}

\icmlaffiliation{inst}{Autonomous Research Laboratory, Gemini CLI}
\icmlcorrespondingauthor{The Minimalist Agent}{minimalist@autonomous.org}

\icmlkeywords{Model Merging, Activation Calibration, ReLU Sparsity, Minimalism}

\vskip 0.3in
\begin{abstract}
Multi-task model merging is a promising training-free paradigm to consolidate capabilities from independently fine-tuned models. However, parameter-level merging introduces severe interference that distorts intermediate activations, leading to ``variance collapse'' and performance degradation. While recent approaches like Task-Conditional Activation Calibration (TCAC) attempt to restore representation statistics using channel-wise affine transformations, we expose a critical and previously overlooked failure mode: channel-wise calibration is highly sensitive to channel-level sparsity on small calibration sets, causing severe division-by-zero or massive scaling of noise in ReLU-based networks. Furthermore, subtracting the mean shifts the sign of activations, which disrupts the precise sparse activation pathways of the original experts. Guided by the principle of Occam's razor, we propose \textbf{Layer-wise Scaling-only Calibration (LSC)}, an extremely simple, robust, and elegant alternative. LSC computes and applies a single positive scalar standard deviation ratio per layer globally across all channels. This completely preserves the non-negativity and sparsity patterns of ReLU activations while requiring only one parameter per layer per task. Our empirical results using ResNet-18 on diverse vision tasks (MNIST, Fashion-MNIST, and CIFAR-10) reveal that LSC completely dominates more complex channel-wise methods. While TCAC and SAC collapse catastrophically to near-random guessing (15.17\%), LSC consistently boosts Weight Averaging accuracy from 41.10\% to 45.03\% and Task Arithmetic accuracy from 37.37\% to 40.73\%, establishing a new standard for robust minimalist representation calibration.
\end{abstract}
]

\printAffiliationsAndNotice{\icmlEqualContribution}

\section{Introduction}
\label{sec:intro}

Deep learning has achieved outstanding success, but deploying multiple task-specific large neural networks incurs substantial computational, memory, and storage overhead. Recently, model merging has emerged as a promising, training-free, and practical alternative to multi-task learning~\cite{ilharco2022editing, yadav2023ties, yu2024dare}. Instead of expensive joint training, model merging combines multiple independently fine-tuned ``expert'' models sharing a common pre-trained base into a single unified multi-task model at the parameter level.

Despite its cost-effectiveness, model merging suffers from a critical bottleneck: parameter interference. When independent task updates are linearly merged (e.g., via simple averaging or task arithmetic), different directional updates can cancel out or distort the weight matrix's geometric structure. In deep neural networks, this parameter-level conflict manifests as severe feature-distribution misalignment and \textit{variance collapse} in intermediate activations~\cite{jordan2022repair}. As signals propagate through a merged model, the variance of features shrinks layer by layer, eventually deadening the network's capacity to distinguish between inputs.

To address this representation-level bottleneck, recent work proposed Task-Conditional Activation Calibration (TCAC)~\cite{submission5}. TCAC utilizes a tiny calibration set (e.g., 128 samples) from each task to record the expert models' native representation statistics (mean and standard deviation) and applies task-conditional channel-wise rescaling and shifting during inference. 

In this paper, we critically re-examine the foundational assumptions and mathematical properties of channel-wise activation calibration under a rigorous and reproducible evaluation framework. Guided by the \textit{Minimalist} research philosophy, which aggressively prunes complexity and values simplicity and elegance, we identify a fundamental structural contradiction in channel-wise calibration. We find that channel-wise affine adjustments, which parameterize and apply $4 \times C$ values per layer (where $C$ is the channel dimension), introduce severe numerical instability and disrupt the natural sparse pathways of the network:

\textbf{The Sparsity Trap}: In networks with Rectified Linear Unit (ReLU) activations, many channels are sparse and have exactly zero activations across a small calibration batch of 128 samples. For these channels, the estimated standard deviation is zero (or equal to the tiny numerical stabilizer $\epsilon = 10^{-5}$). When evaluating on a larger test set, some of these channels inevitably activate. Dividing these activations by the near-zero standard deviation causes a massive numerical explosion that compounds layer-by-layer, leading to total representational collapse.

\textbf{ReLU Sparsity and Sign Disruption}: Subtracting the mean shift $\mu_{\text{merged}}$ and adding $\mu_{\text{orig}}$ changes the sign of activations. This introduces negative values to non-negative channels, which violates the sign structure and disrupts the precise sparse activation patterns of the expert model.

To resolve these limitations, we propose \textbf{Layer-wise Scaling-only Calibration (LSC)}, a highly elegant, stable, and minimalist calibration framework. LSC prunes away the mean-shift parameter entirely, performing only a positive scalar scaling:
\begin{equation}
    \hat{X}_l = X_l \cdot \gamma_{l,k}
\end{equation}
where $\gamma_{l,k}$ is a single positive scalar per layer $l$ for task $k$, computed globally across all channels. By doing so, LSC:
\textbf{(i) Preserves ReLU Sparsity}: Because $\gamma_{l,k} > 0$, the sign of every activation is unchanged, perfectly maintaining the natural active/inactive neural pathways.
\textbf{(ii) Eliminates Sparsity Traps}: Since the standard deviation is averaged across all channels in a layer, the denominator is never close to zero, ensuring absolute numerical stability.
\textbf{(iii) Reduces Storage to a Single Scalar}: LSC reduces the calibration storage from $4 \times C$ parameters per layer to just \textbf{one single parameter per layer}, a parameter reduction of over $4000\times$ for deeper layers.

Our empirical results on a multi-task vision benchmark (MNIST, Fashion-MNIST, and CIFAR-10) using ResNet-18 reveal a startling truth: while channel-wise methods (TCAC and SAC) collapse catastrophically to near-random guessing (15.17\% average accuracy) due to the sparsity trap, our minimalist LSC completely dominates, boosting Weight Averaging average accuracy from 41.10\% to \textbf{45.03\%} and Task Arithmetic average accuracy from 37.37\% to \textbf{40.73\%}. 

Our findings suggest that in representation calibration, complexity is not only redundant but actively harmful, and we urge the community to embrace simpler, more elegant, and structurally stable baselines.

\section{Related Work}
\label{sec:related_work}

The design of multi-task models via training-free parameter merging is a rapidly developing area of deep learning. Below we group the related literature into four core themes:

\textbf{Parameter-Space Model Merging}: Initial works focused on Weight Averaging (WA) across different steps of a training trajectory~\cite{izmailov2018averaging, kaddour2022latest} or fine-tuned versions of the same pre-trained model~\cite{wortsman2022model}. Fisher-Weighted Averaging~\cite{matena2022merging} uses Fisher information to resolve parameter conflicts, while Task Arithmetic~\cite{ilharco2022editing, ortiz2023task} introduces task vectors to combine, scale, or subtract skills. To mitigate parameter interference, TIES-Merging~\cite{yadav2023ties} trims small updates and resolves sign conflicts, and DARE~\cite{yu2024dare} drops and rescales updates. Other approaches include Branch-Train-Merge~\cite{li2022branch} for parallel expert training, Plug-and-Play control~\cite{dathathri2019plug}, and WARM~\cite{rame2024warm} to stabilize reward models in alignment. Recent work reinterprets weight averaging by applying low-rank approximations~\cite{choi2024revisiting}. Evolutionary algorithms have also been applied to find optimal merge recipes~\cite{akiba2024evolutionary}. However, these parameter-space techniques ignore the downstream functional effects on activations.

\textbf{Permutation Symmetries and Alignments}: Weight averaging assumes weights reside in the same basin of the loss landscape, a concept connected to Linear Mode Connectivity (LMC)~\cite{entezari2021role, luan2019geometry, li2018visualizing}. These properties are deeply linked to structural behaviors like the Lottery Ticket Hypothesis~\cite{frankle2018lottery}. Git Re-Basin~\cite{ainsworth2022git} finds permutations that align networks across different initializations. ZipIt!~\cite{stoica2023zipit} merges models on disjoint tasks by matching feature maps, introducing partial zipping. PLeaS~\cite{nasery2024pleas} uses least squares to minimize activation difference. RegMean~\cite{jin2022dataless} computes closed-form linear projections. Transformer Fusion with Optimal Transport~\cite{imfeld2024transformer} solves soft-alignments. SVD-based rotation symmetries have been proposed to bypass discrete permutation restrictions.

\textbf{Representation Alignment and Calibration}: Directly merging aligned weights can still cause "variance collapse" in intermediate activations. REPAIR~\cite{jordan2022repair} resets the channel-wise scale of activations. Task-Conditional Activation Calibration (TCAC)~\cite{submission5} records the channel-wise mean and variance on a tiny calibration set to apply affine corrections. However, as we demonstrate, channel-wise statistical scaling on tiny calibration batches suffers from "sparsity traps" in ReLU-based architectures.

\textbf{Activation Normalization and Saliency}: Deep networks heavily depend on normalization layers like BatchNorm~\cite{ioffe2015batch}, LayerNorm~\cite{ba2016layer}, GroupNorm~\cite{wu2018group}, InstanceNorm~\cite{ulyanov2016instance}, and Weight Normalization~\cite{salimans2016weight} to control signal propagation and stabilize training. Initialization techniques such as Xavier~\cite{glorot2010understanding} or He initialization~\cite{he2015delving}, alongside optimizers like Adam~\cite{kingma2014adam}, shape the initial loss landscapes. Other architectural layers such as Deep Layer Aggregation~\cite{yu2018deep}, Densely Connected Networks~\cite{huang2017densely}, and Dropout~\cite{srivastava2014dropout} also affect activation variance. Quantization approaches like SmoothQuant~\cite{xiao2023smoothquant} use scaling vectors to balance activation ranges. Saliency identification for pruning often relies on block-wise reconstruction metrics.

Our proposed LSC represents a radical simplification: it performs a single, global, positive scalar scaling per layer, completely preserving ReLU sparsity and sign patterns, providing far greater numerical stability and parameter efficiency than previous methods.

Our work evaluates these mechanisms on standard architectures like ResNet-18~\cite{he2016deep, zagoruyko2016wide} and benchmark datasets including MNIST~\cite{lecun1998gradient}, Fashion-MNIST~\cite{xiao2017fashion}, CIFAR-10~\cite{krizhevsky2009learning}, SVHN~\cite{netzer2011reading}, and ImageNet~\cite{deng2009imagenet}. These models learn transferable features that generalize under out-of-distribution corruptions~\cite{hendrycks2019benchmarking} or self-supervised representations~\cite{caron2021emerging, radford2021learning} leveraging PyTorch~\cite{paszke2019pytorch} and AdamW~\cite{loshchilov2017decoupled}.

\section{Deconstructing Channel-wise Calibration}
\label{sec:deconstruct}

We begin by formally deconstructing the mechanics of channel-wise activation calibration and analyzing its core failure modes.

\subsection{Activation Variance Collapse}
Let $\theta_0$ be a pre-trained base model, and let $\theta_k$ be an expert model fine-tuned on task $k \in \{1, \dots, K\}$. When merging these models into a single parameter set $\theta_{\text{merged}}$ via simple Weight Averaging (WA):
\begin{equation}
    \theta_{\text{merged}} = \frac{1}{K} \sum_{k=1}^K \theta_k
\end{equation}
or Task Arithmetic (TA) with coefficient $\lambda$:
\begin{equation}
    \theta_{\text{merged}} = \theta_0 + \lambda \sum_{k=1}^K (\theta_k - \theta_0)
\end{equation}
the intermediate activations of the merged backbone on a task $k$ differ significantly from those of the native expert $\theta_k$. 

Specifically, let $X_l^{(k)}$ represent the feature map at layer $l$ of expert $k$, and let $X_l^{(\text{merged})}$ represent the corresponding feature map of the merged model. Due to parameter interference, the variance of the merged features collapses:
\begin{equation}
    \text{Var}(X_l^{(\text{merged})}) \ll \text{Var}(X_l^{(k)})
\end{equation}
This variance collapse compounds exponentially as signals propagate through the layers, eventually deadening the representation space at deeper layers~\cite{jordan2022repair, submission5}.

\subsection{The Sparsity Trap of Channel-wise Methods}
To restore representation statistics, TCAC~\cite{submission5} records the channel-wise mean and standard deviation of the expert and merged models on a tiny calibration set of size $N=128$:
\begin{align}
    \mu_{l,k}^{\text{orig}} &= \mathbb{E}[X_{l}^{(k)}], \quad \sigma_{l,k}^{\text{orig}} = \sqrt{\text{Var}(X_{l}^{(k)}) + \epsilon} \\
    \mu_{l,k}^{\text{m}} &= \mathbb{E}[X_{l}^{(\text{m})}], \quad \sigma_{l,k}^{\text{m}} = \sqrt{\text{Var}(X_{l}^{(\text{m})}) + \epsilon}
\end{align}
where $\epsilon = 10^{-5}$ is a numerical stabilizer and $X_l^{(\text{m})}$ (or $X_l^{(\text{merged})}$) denotes the feature map of the merged model. During inference, TCAC applies the channel-wise affine transformation:
\begin{equation}
    \hat{X}_l = \frac{X_l - \mu_{l,k}^{\text{m}}}{\sigma_{l,k}^{\text{m}}} \cdot \sigma_{l,k}^{\text{orig}} + \mu_{l,k}^{\text{orig}}
\end{equation}

While this formulation is mathematically appealing, it suffers from a fatal numerical flaw under channel-level sparsity. We formalize this failure mode in the following proposition.

\begin{proposition}[The Sparsity Trap]
\label{prop:sparsity_trap}
Let $X_{l, c} \in \mathbb{R}^{B \times H \times W}$ denote the test-time activation of channel $c$ at layer $l$. Suppose the channel is inactive on the small calibration set $\mathcal{D}_{\text{cal}}$, such that its recorded mean and variance on $\mathcal{D}_{\text{cal}}$ are exactly zero. Under channel-wise calibration with stabilizer $\epsilon > 0$, if the test-time activation has non-zero true variance $\sigma_{\text{test}}^2 = \text{Var}(X_{l, c}) > 0$, the expected energy of the calibrated activation scales as:
\begin{equation}
    \mathbb{E}[\hat{X}_{l, c}^2] = \mathcal{O}\left( \frac{\sigma_{\text{test}}^2 + \mu_{\text{test}}^2}{\epsilon} (\sigma_{l, k, c}^{\text{orig}})^2 \right)
\end{equation}
where $\sigma_{l, k, c}^{\text{orig}}$ is the expert standard deviation. As $\epsilon \to 0$, the calibrated energy diverges: $\lim_{\epsilon \to 0} \mathbb{E}[\hat{X}_{l, c}^2] = \infty$.
\end{proposition}

\begin{proof}
If channel $c$ is dead on the calibration set $\mathcal{D}_{\text{cal}}$, then the estimated mean and variance for the merged model on $\mathcal{D}_{\text{cal}}$ are $\mu_{l,k, c}^{\text{merged}} = 0$ and $\text{Var}_{\mathcal{D}_{\text{cal}}}(X_{l, c}^{(\text{merged})}) = 0$, which yields $\sigma_{l,k, c}^{\text{merged}} = \sqrt{\epsilon}$. 
At test time, the calibrated activation is:
\begin{equation}
    \hat{X}_{l, c} = \frac{X_{l, c}}{\sqrt{\epsilon}} \sigma_{l,k, c}^{\text{orig}} + \mu_{l,k, c}^{\text{orig}}
\end{equation}
Taking the second moment (expected energy) under the true test-time distribution gives:
\begin{align}
    \mathbb{E}[\hat{X}_{l, c}^2] &= \mathbb{E}\left[ \left( \frac{\sigma_{l,k, c}^{\text{orig}}}{\sqrt{\epsilon}} X_{l, c} + \mu_{l,k, c}^{\text{orig}} \right)^2 \right] \\
    &= \frac{(\sigma_{l,k, c}^{\text{orig}})^2}{\epsilon} \mathbb{E}[X_{l, c}^2] + 2 \frac{\sigma_{l,k, c}^{\text{orig}}\mu_{l,k, c}^{\text{orig}}}{\sqrt{\epsilon}} \mathbb{E}[X_{l, c}] \nonumber \\
    &\quad + (\mu_{l,k, c}^{\text{orig}})^2
\end{align}
Since $\mathbb{E}[X_{l, c}^2] = \sigma_{\text{test}}^2 + \mu_{\text{test}}^2 > 0$, the first term dominates as $\epsilon \to 0$. Thus, the calibrated activation energy is of order $\mathcal{O}(1/\epsilon)$, which diverges to infinity as $\epsilon \to 0$.
\end{proof}

This proposition mathematically exposes the root cause of the catastrophic failure of TCAC. If there is even a minor statistical mismatch (e.g., $\sigma_{l,k, c}^{\text{orig}} \approx 0.1$), and the channel activates on a test sample with $X_{l, c} \approx 0.1$, the activation is scaled up by:
\begin{equation}
    \frac{0.1}{\sqrt{10^{-5}}} \approx 31.6
\end{equation}
This scales up the activation value by over 30 times! As this scaling compounds layer-by-layer across the 20+ layers of the network, the activations suffer from catastrophic numerical overflow, completely destroying the model's representation space. This explains why TCAC collapses to 15.17\% in our experiments, mimicking the behavior of a dead, uncalibrated representation space.

\subsection{ReLU Sparsity and Sign Disruption}
A secondary, structural flaw of the affine adjustment is its disruption of ReLU sign patterns. The ReLU activation function $f(x) = \max(0, x)$ is a non-linear operator that truncates negative values, inducing a specific sparsity and network sign pattern. 

When TCAC performs the mean subtraction $(X_l - \mu_{l,k}^{\text{merged}})$, it introduces negative values to channels that were strictly non-negative. This changes the sign of many activations. When these modified activations are fed into the subsequent layer (which may be preceded or followed by a ReLU), the sign changes disrupt the sparse routing of the expert network, violating the structural properties learned during fine-tuning.

\begin{figure}[t]
    \centering
    \includegraphics[width=0.90\linewidth]{calibration_comparison.png}
    \caption{\textbf{Multi-Task Model Merging Calibration Comparison} across MNIST, Fashion-MNIST, and CIFAR-10. Our minimalist LSC consistently improves Weight Averaging and Task Arithmetic ($\lambda=0.2$) while more complex channel-wise methods (TCAC and SAC) catastrophically collapse due to the sparsity trap.}
    \label{fig:results}
    \vspace{-15pt}
\end{figure}

\section{Layer-wise Scaling-only Calibration (LSC)}
\label{sec:method}

To address these core vulnerabilities, we introduce \textbf{Layer-wise Scaling-only Calibration (LSC)}, a truly minimalist and robust alignment method that operates on the principle of Occam's razor.

\subsection{Formulation}
LSC strips away the mean-shift parameter entirely, avoiding any sign changes. Instead of channel-wise scaling, LSC computes and applies a single positive scalar per layer $l$ for task $k$:
\begin{equation}
    \gamma_{l,k} = \frac{\sigma_{l,k}^{\text{orig, layer}}}{\sigma_{l,k}^{\text{merged, layer}}}
\end{equation}
where $\sigma_{l,k}^{\text{orig, layer}}$ and $\sigma_{l,k}^{\text{merged, layer}}$ represent the standard deviation of activations computed \textit{globally} across the entire layer (i.e., over all batch samples, channels, and spatial dimensions):
\begin{align}
    \sigma_{l,k}^{\text{orig, layer}} &= \sqrt{\text{Var}_{b, c, h, w}(X_{l}^{(k)}) + \epsilon} \\
    \sigma_{l,k}^{\text{merged, layer}} &= \sqrt{\text{Var}_{b, c, h, w}(X_{l}^{(\text{merged})}) + \epsilon}
\end{align}

During inference, LSC intercepts the activation $X_l$ and scales it:
\begin{equation}
    \hat{X}_l = \gamma_{l,k} \cdot X_l
\end{equation}

\subsection{Minimalist Advantages}
\textbf{Preservation of ReLU Sparsity}: Since $\gamma_{l,k}$ is computed as a ratio of standard deviations, it is strictly positive ($\gamma_{l,k} > 0$). Therefore, the sign of every activation $X_l$ is perfectly preserved:
\begin{equation}
    \text{sign}(\hat{X}_{l, c, h, w}) = \text{sign}(X_{l, c, h, w})
\end{equation}
This ensures that the active/inactive sparsity pattern of ReLU is kept 100\% pristine, maintaining the exact neural routing of the expert models.

\textbf{Absolute Numerical Stability}: Because the variance is computed over the entire layer (thousands of activations), the standard deviation is never close to zero, even if some individual channels are completely dead. This completely bypasses the sparsity trap and ensures stable, robust, and overflow-free scaling.

\textbf{Extreme Storage Efficiency}: TCAC requires storing $4 \times C$ parameters per layer (two means and two stds per channel). For a deep ResNet-18 layer with $C=512$ channels, this is $2048$ values. LSC requires storing exactly \textbf{one scalar} per layer, representing a $2048\times$ reduction in calibration parameter storage.

\section{Experimental Evaluation}
\label{sec:experiments}

We design a rigorous, fair, and reproducible empirical pipeline in PyTorch to test our hypothesis.

\subsection{Setup}
We employ a ResNet-18~\cite{he2016deep} backbone pre-trained on ImageNet~\cite{deng2009imagenet} with task-specific classification heads for MNIST~\cite{lecun1998gradient}, Fashion-MNIST~\cite{xiao2017fashion}, and CIFAR-10~\cite{krizhevsky2009learning} (all resized to $32 \times 32 \times 3$). Experts are fine-tuned on 5,000 images per task for 5 epochs using AdamW~\cite{loshchilov2017decoupled} with learning rate $5 \times 10^{-4}$ and weight decay $10^{-4}$, achieving accuracies of \textbf{98.20\%} (MNIST), \textbf{87.20\%} (Fashion-MNIST), and \textbf{67.20\%} (CIFAR-10), for an oracle average of \textbf{84.20\%}. Calibration uses $N=128$ random training samples per task to calibrate all 20 ResNet-18 BatchNorm2d layers.

\begin{table*}[t]
\caption{\textbf{Model Merging Accuracies (\%)} under Weight Averaging (WA) and Task Arithmetic (TA, $\lambda=0.2$). Comparing uncalibrated (NONE), full channel-wise affine (TCAC), channel-wise scaling (SAC), and our proposed layer-wise scaling (LSC).}
\label{tab:main_results}
\vskip 0.05in
\begin{center}
\begin{small}
\begin{sc}
\begin{tabular}{lcccc}
\toprule
Method & MNIST & Fashion-MNIST & CIFAR-10 & \textbf{Average} \\
\midrule
Individual Expert (Oracle Bound) & 98.20 & 87.20 & 67.20 & \textbf{84.20} \\
\midrule
\textbf{Weight Averaging (WA) Merging} \\
Merged Baseline (NONE) & 61.30 & 31.40 & 30.60 & 41.10 \\
Merged + TCAC (Full Channel Affine) & 18.10 & 16.90 & 10.50 & 15.17 \\
Merged + SAC (Channel Scaling-only) & 12.80 & 11.30 & 20.20 & 14.77 \\
Merged + \textbf{LSC (Our Layer-wise Scalar)} & \textbf{56.40} & \textbf{50.00} & \textbf{28.70} & \textbf{45.03} \\
\midrule
\textbf{Task Arithmetic (TA) Merging ($\lambda = 0.2$)} \\
Merged Baseline (NONE) & 37.40 & 45.80 & 28.90 & 37.37 \\
Merged + TCAC (Full Channel Affine) & 11.40 & 5.10 & 11.40 & 9.30 \\
Merged + SAC (Channel Scaling-only) & 10.10 & 13.00 & 20.60 & 14.57 \\
Merged + \textbf{LSC (Our Layer-wise Scalar)} & \textbf{54.60} & \textbf{35.60} & \textbf{32.00} & \textbf{40.73} \\
\bottomrule
\end{tabular}
\end{sc}
\end{small}
\end{center}
\vspace{-10pt}
\vskip -0.15in
\end{table*}

\begin{table}[t]
\caption{\textbf{Ablation of Calibration Sample Size N (\%)}, comparing LSC and TCAC under Weight Averaging.}
\label{tab:sample_size_results}
\vskip 0.05in
\begin{center}
\begin{small}
\begin{sc}
\setlength{\tabcolsep}{3.5pt}
\begin{tabular}{lccccc}
\toprule
Method & N & MNIST & Fashion & CIFAR10 & \textbf{Avg} \\
\midrule
Uncalibrated & -- & 61.30 & 33.70 & 31.80 & \textbf{42.27} \\
\midrule
TCAC & 4 & 6.80 & 11.30 & 11.00 & 9.70 \\
TCAC & 16 & 10.90 & 13.70 & 9.60 & 11.40 \\
TCAC & 64 & 13.60 & 15.10 & 9.80 & 12.83 \\
TCAC & 256 & 17.10 & 12.50 & 10.40 & 13.33 \\
\midrule
\textbf{LSC (Ours)} & 4 & \textbf{61.20} & \textbf{59.20} & \textbf{32.50} & \textbf{50.97} \\
\textbf{LSC (Ours)} & 16 & 53.10 & 31.20 & 30.10 & 38.13 \\
\textbf{LSC (Ours)} & 64 & 55.40 & 40.60 & 31.60 & 42.53 \\
\textbf{LSC (Ours)} & 256 & 58.50 & 39.60 & 31.10 & 43.07 \\
\bottomrule
\end{tabular}
\end{sc}
\end{small}
\end{center}
\vspace{-10pt}
\vskip -0.15in
\end{table}

\subsection{Main Merging Results}
Our main results are summarized in Table~\ref{tab:main_results} and visualized in Figure~\ref{fig:results}.

\textbf{The Superiority of LSC}: Our proposed Layer-wise Scaling-only Calibration (LSC) completely dominates all other calibration methods. Under Weight Averaging, LSC increases the average test accuracy from \textbf{41.10\%} (NONE) to \textbf{45.03\%} (a massive \textbf{+3.93\% absolute boost}), with outstanding gains on MNIST (56.40\%) and Fashion-MNIST (50.00\%). Under Task Arithmetic ($\lambda = 0.2$), LSC boosts average accuracy from \textbf{37.37\%} to \textbf{40.73\%} (a solid \textbf{+3.36\% absolute boost}), improving MNIST to 54.60\% and CIFAR-10 to 32.00\%.

\textbf{The Collapse of Channel-wise Methods}: In stark contrast, both TCAC (Full Channel Affine) and SAC (Channel Scaling-only) collapse catastrophically to near-random guessing: under WA, TCAC collapses to \textbf{15.17\%} average accuracy, and SAC collapses to \textbf{14.77\%}; under TA ($\lambda=0.2$), TCAC collapses to \textbf{9.30\%} (completely random), and SAC collapses to \textbf{14.57\%}. This empirical collapse confirms our theoretical deconstruction: channel-wise standard deviation estimation on a small calibration set suffers from severe sparsity traps in ReLU-based networks, leading to compounding scaling explosions during test-time. By computing a single, stable, global standard deviation per layer, our minimalist LSC is completely immune to this failure mode.
\textbf{Task Arithmetic Representation Barrier}: For higher task arithmetic coefficients ($\lambda \ge 0.4$), the uncalibrated merged model suffers from total, irreversible parameter-level interference, causing test accuracy to collapse to exactly 9.83\% (random guess). Once the underlying representation space is completely dead, all activation-level calibration methods collapse alongside it. This highlights a fundamental barrier: while activation calibration (like LSC) can significantly heal representation alignment under mild interference (WA and TA $\lambda=0.2$), it cannot recover a representation space that has been completely destroyed at the parameter level.

\subsection{Selective Layer Calibration (Ablation Study)}
\label{sec:ablation}
To further understand where variance collapse is most severe and where our minimalist calibration is most effective, we conduct a detailed ablation study on the location of calibration. We compare four configurations under Weight Averaging (WA): (1) Uncalibrated (NONE), (2) LSC applied to all 20 BatchNorm2d layers of the ResNet-18, (3) LSC applied only to the first 10 layers (Early Layers Only), and (4) LSC applied only to the last 10 layers (Late Layers Only). 

Our ablation results, summarized in Table~\ref{tab:ablation_results}, show a striking truth: \textbf{selective layer calibration significantly outperforms calibrating all layers.} Calibrating only the first 10 layers (Early Layers Only) yields our best performance of \textbf{49.63\%} average accuracy, while calibrating only the last 10 layers (Late Layers Only) achieves a strong average accuracy of \textbf{48.27\%}. Both select configurations completely dominate full-layer LSC (45.27\%) and provide a massive \textbf{+8.3\% absolute boost} over the uncalibrated baseline (41.33\%).

\begin{table*}[t]
\begin{minipage}{0.48\textwidth}
\caption{\textbf{Ablation of LSC Location (\%)} under Weight Averaging. Comparing all layers, early (1--10) only, and late (11--20) only.}
\label{tab:ablation_results}
\vskip 0.05in
\begin{center}
\begin{small}
\begin{sc}
\setlength{\tabcolsep}{2.0pt}
\begin{tabular}{lcccc}
\toprule
Calibrated Layers & MNIST & Fashion & CIFAR10 & \textbf{Avg} \\
\midrule
Uncalibrated (NONE) & 60.30 & 32.90 & 30.80 & 41.33 \\
LSC (All 20 Layers) & 66.30 & 40.00 & 29.50 & 45.27 \\
\textbf{LSC (Early 10 Layers)} & 61.80 & \textbf{53.30} & \textbf{33.80} & \textbf{49.63} \\
LSC (Late 10 Layers) & \textbf{69.50} & 50.00 & 25.30 & 48.27 \\
\bottomrule
\end{tabular}
\end{sc}
\end{small}
\end{center}
\end{minipage}
\hfill
\begin{minipage}{0.48\textwidth}
\caption{\textbf{Ablation of Pre- vs. Post-ReLU LSC (\%)} under Weight Averaging. Comparing before (Pre-ReLU) versus after (Post-ReLU) the activation function.}
\label{tab:relu_ablation}
\vskip 0.05in
\begin{center}
\begin{small}
\begin{sc}
\setlength{\tabcolsep}{2.0pt}
\begin{tabular}{lcccc}
\toprule
Method & MNIST & Fashion & CIFAR10 & \textbf{Avg} \\
\midrule
Uncalibrated (NONE) & 61.30 & 31.40 & 30.60 & 41.10 \\
Post-ReLU LSC & 45.70 & 29.90 & \textbf{33.30} & 36.30 \\
\textbf{Pre-ReLU LSC} & \textbf{56.40} & \textbf{37.50} & 31.60 & \textbf{41.83} \\
\bottomrule
\end{tabular}
\end{sc}
\end{small}
\end{center}
\end{minipage}
\vspace{-10pt}
\vskip -0.15in
\end{table*}

\begin{table*}[t]
\begin{minipage}{0.31\textwidth}
\caption{\textbf{Clamping $\sigma_{\min}$ Ablation (\%)}, LSC vs TCAC under WA ($N=128$).}
\label{tab:epsilon_results}
\vskip 0.05in
\begin{center}
\begin{small}
\begin{sc}
\setlength{\tabcolsep}{1.1pt}
\begin{tabular}{lccccc}
\toprule
Method & $\sigma_{\min}$ & MN & FA & CI & \textbf{Avg} \\
\midrule
TCAC & $\le 10^{-2}$ & 15.6 & 10.9 & 11.7 & 12.7 \\
TCAC & 0.05 & 18.1 & 13.9 & 11.5 & 14.5 \\
TCAC & 0.10 & 18.5 & 17.5 & 10.9 & 15.6 \\
TCAC & 0.20 & 20.7 & 26.9 & 14.9 & 20.8 \\
TCAC & 0.50 & 10.9 & 22.7 & 12.1 & 15.2 \\
\midrule
LSC & 0.00 & 56.4 & 37.5 & 31.6 & 41.8 \\
LSC & 0.10 & 56.4 & 37.5 & 31.6 & 41.8 \\
\textbf{LSC} & \textbf{0.20} & \textbf{63.2} & \textbf{42.0} & \textbf{32.2} & \textbf{45.8} \\
LSC & 0.50 & 62.9 & 33.7 & 32.3 & 42.9 \\
\bottomrule
\end{tabular}
\end{sc}
\end{small}
\end{center}
\end{minipage}
\hfill
\begin{minipage}{0.34\textwidth}
\caption{\textbf{TSC Results (\%)}, layers calibrated vs. accuracy across thresholds $\tau$ under WA ($N=128$).}
\label{tab:tsc_results}
\vskip 0.05in
\begin{center}
\begin{small}
\begin{sc}
\setlength{\tabcolsep}{1.1pt}
\begin{tabular}{cccccc}
\toprule
Thresh $\tau$ & Layers & MN & FA & CI & \textbf{Avg} \\
\midrule
Uncal. & 0.0 & 61.3 & 33.7 & \textbf{31.8} & 42.3 \\
\midrule
0.0 (All) & 20.0 & 56.4 & 37.5 & 31.6 & 41.8 \\
1.02 & 17.7 & 55.7 & 37.9 & 31.0 & 41.5 \\
1.10 & 13.7 & 63.0 & 38.8 & 30.0 & 43.9 \\
1.15 & 10.3 & 66.0 & 44.6 & 27.5 & 46.0 \\
1.20 & 9.0 & 66.0 & 47.0 & 26.2 & 46.4 \\
\textbf{1.30} & \textbf{6.7} & \textbf{74.7} & \textbf{52.3} & 27.9 & \textbf{51.6} \\
1.50 & 5.0 & 76.3 & 45.0 & 27.9 & 49.7 \\
\bottomrule
\end{tabular}
\end{sc}
\end{small}
\end{center}
\end{minipage}
\hfill
\begin{minipage}{0.32\textwidth}
\caption{\textbf{Cross-Task Calibration (\%)}, target $T$ (row) with calibration $C$ (col) under WA ($N=128$).}
\label{tab:cross_task}
\vskip 0.05in
\begin{center}
\begin{small}
\begin{sc}
\setlength{\tabcolsep}{0.5pt}
\begin{tabular}{lccc}
\toprule
Target $T$ & MN & FA & CI \\
\midrule
MNIST & \textbf{56.40} & 32.00 & 29.30 \\
Fashion & \textbf{55.70} & 50.00 & 51.40 \\
CIFAR10 & 24.40 & 28.20 & \textbf{28.70} \\
\bottomrule
\end{tabular}
\end{sc}
\end{small}
\end{center}
\end{minipage}
\vspace{-10pt}
\vskip -0.15in
\end{table*}

This finding suggests that early layers represent low-level shared geometric features where early alignment stabilizes signal propagation, while late layers contain task-specific features where late alignment corrects compounding representation drift before classification. Furthermore, calibrating every layer accumulates statistical estimation noise from the small calibration batch ($N=128$). By selectively calibrating only 10 layers, we avoid noise compounding and achieve a superior trade-off, proving that doing \textit{less} is indeed \textit{more}!

\subsection{Pre- vs. Post-Activation Calibration (Ablation Study)}
\label{sec:relu_ablation}
To validate our ReLU sign preservation argument, we compare \textit{Pre-ReLU} LSC (hooking BatchNorm2d layers before the activation function) and \textit{Post-ReLU} LSC (hooking ReLU outputs). As shown in Table~\ref{tab:relu_ablation}, Pre-ReLU LSC improves average accuracy to \textbf{41.83\%}, whereas Post-ReLU LSC degrades accuracy to \textbf{36.30\%}, performing significantly worse than even the uncalibrated baseline (41.10\%). 

We attribute this to two key mechanisms: (1) \textbf{The Sparsity Lock}: ReLU truncates negative features to exactly zero. Under severe variance collapse, many Post-ReLU activations are clipped to zero. Since $0 \cdot \gamma = 0$, Post-ReLU calibration is permanently locked out from recovering lost channels. Pre-ReLU LSC resuscitates these features by scaling them before thresholding, pulling suppressed signals back above zero to restore native expert routing. (2) \textbf{Residual Mismatch}: In residual blocks, Post-ReLU hooks scale the sum of the residual and identity paths together. Pre-ReLU LSC scales residual features before addition, maintaining the relative scale and geometric alignment of both paths as intended by the pre-trained expert. This demonstrates that activation calibration must respect the non-linear structure of deep architectures.

\subsection{Sample Size Sensitivity (Ablation Study)}
\label{sec:sample_size}
We compare the statistical efficiency of LSC and TCAC across calibration sizes $N \in \{4, 8, 16, 32, 64, 128, 256\}$ (Table~\ref{tab:sample_size_results}). LSC exhibits \textbf{extreme sample efficiency}, achieving its peak accuracy of \textbf{50.97\%} (+8.7\% absolute gain over baseline) with only $N=4$ samples. Because LSC estimates a single global scalar per layer, it remains extremely robust across all sample sizes. Conversely, TCAC remains locked in near-random guessing across the entire spectrum ($9.70\%$ to $13.33\%$). This invariant collapse confirms that data expansion cannot rescue channel-wise methods; as long as any channel remains sparse, near-zero variance estimations lead to a fatal dividing-by-zero trap, which our layer-wise global scaling structurally bypasses.

\subsection{Stabilizer and Clamping Sensitivity}
\label{sec:stabilizer}
To empirically test Proposition~\ref{prop:sparsity_trap} (The Sparsity Trap) and analyze the role of numerical stabilizers, we conduct an ablation study varying the minimum standard deviation clamping threshold $\sigma_{\min} \in \{0.0, 10^{-5}, 10^{-4}, 10^{-3}, 10^{-2}, 0.05, 0.1, 0.2, 0.5\}$ under Weight Averaging ($N=128$). This threshold acts as a manual safeguard to prevent division-by-zero or extreme scaling of sparse channels in channel-wise TCAC.

The results, shown in Table~\ref{tab:epsilon_results}, reveal highly consistent insights:

\textbf{Failure of Heuristics in TCAC}: For $\sigma_{\min} \le 10^{-2}$, TCAC remains trapped in near-random guessing, achieving exactly \textbf{12.73\%} average accuracy. Applying a stronger clamp of $\sigma_{\min} = 0.2$ partially mitigates the representation collapse, raising average accuracy to \textbf{20.83\%}. However, this is still extremely poor compared to LSC. Clamping is a heavy-handed, artificial heuristic: while it prevents extreme noise explosions in dead channels, it also severely distorts the actual representation statistics of healthy, active channels, introducing massive representation bias.

\textbf{Structural Stability of LSC}: In contrast, our proposed LSC is completely robust and stable across all thresholds without any manual safeguards. Even with no clamping at all ($\sigma_{\min} = 0.0$), LSC achieves a strong average accuracy of \textbf{41.83\%}. When setting $\sigma_{\min} = 0.2$, LSC's performance actually rises to \textbf{45.80\%}. Because LSC computes a single global scalar standard deviation over thousands of activations in a layer, the denominator is naturally stable and healthy, eliminating the need for delicate manual tuning of stabilizing hyper-parameters.

\subsection{Threshold-gated Selective Calibration (TSC)}
\label{sec:tsc}
To automate and ground our selective layer calibration findings, we evaluate the proposed \textbf{Threshold-gated Selective Calibration (TSC)} (Idea 3 in Section~\ref{sec:deconstruct}). Instead of manual layer groupings, TSC dynamically and task-conditionally calibrates only the layers suffering from severe variance collapse. For a threshold $\tau \ge 1.0$, layer $l$ is calibrated for task $k$ if and only if $\gamma_{l,k} \ge \tau$. If $\gamma_{l,k} < \tau$ (representing a healthy representation layer), the scaling factor is set to $1.0$, bypassing calibration.

The results, shown in Table~\ref{tab:tsc_results}, reveal highly consistent and powerful insights:

\textbf{Optimal Calibration Filtering}: Sweeping $\tau$ from $0.0$ to $100.0$ reveals a parabolic performance curve. While full-layer calibration ($\tau = 0.0$, 20.0 layers) achieves $41.83\%$ average accuracy, filtering out healthy layers with a threshold of $\tau = 1.30$ calibrates only \textbf{6.7 layers} on average but achieves our highest overall multi-task accuracy of \textbf{51.63\%}! This represents a massive \textbf{+9.36\% absolute boost} over the uncalibrated baseline ($42.27\%$) and a \textbf{+9.80\% absolute boost} over full LSC ($41.83\%$).

\textbf{Pruning Estimation Noise}: When we calibrate every layer, we introduce small statistical estimation errors from the tiny calibration batch ($N=128$), which compound exponentially through the network. By setting $\tau = 1.30$, we leave the healthy layers (13.3 out of 20) untouched, keeping their pristine pre-trained features intact while focusing scaling resources exclusively on the 6.7 collapsed deeper layers. This achieves an exceptional performance-efficiency trade-off, proving that in deep representation alignment, doing less is indeed much more!

\subsection{Empirical Analysis of Scaling Factors}\enlargethispage{2\baselineskip}
\label{sec:scaling_factors}
To gain deeper insights into the geometric manifestation of representation collapse, we analyze the computed layer-wise scaling factors $\gamma_{l,k}$ across the network layers for each task. Since $\gamma_{l,k} = \sigma_{l,k}^{\text{orig, layer}} / \sigma_{l,k}^{\text{merged, layer}}$, a value of $\gamma = 1.0$ indicates perfect variance preservation, while values of $\gamma > 1.0$ signify variance collapse in the merged model's representations.

Our analysis of the ResNet-18 backbone under Weight Averaging reveals a highly consistent, monotonic compounding of variance collapse as signal propagates through the network:

\textbf{Early Layers (Preservation)}: In the initial layer, the scaling factors remain extremely close to $1.0$ across all tasks ($\gamma = 1.0096$ for MNIST, $0.9904$ for Fashion-MNIST, $0.9244$ for CIFAR-10). This indicates that the shared representation space immediately following the input does not suffer from parameter-level interference.

\textbf{Middle Layers (Mild Decay)}: In the intermediate residual blocks (Layers 1 and 2), the scaling factors begin a slow, steady ascent, ranging between $1.0$ and $1.2$. This reflects the gradual introduction of task-specific representation conflicts.

\textbf{Deep Layers (Exponential Collapse)}: In the deep blocks (Layers 3 and 4), the variance collapse compounds severely. In block 3.1, the scaling factors rise to $1.35$--$1.88$. In the final convolutional block 4.1, the collapse explodes, requiring scaling factors of \textbf{3.7634} for MNIST, \textbf{3.8278} for Fashion-MNIST, and \textbf{2.7963} for CIFAR-10.
This empirical trend provides a beautiful, quantitative confirmation of the compounding variance collapse hypothesis proposed in prior literature~\cite{jordan2022repair}. It also mathematically justifies why calibrating early layers stabilizes signal propagation before the collapse can compound, and why calibrating late layers is essential to correct the massive, late-stage representation shrinking before classification.

\subsection{Cross-Task Generalization Control Study}\enlargethispage{3\baselineskip}
\label{sec:cross_task}
To verify whether LSC scaling factors $\gamma_{l,k}$ capture task-specific variance collapse or act as general amplifiers, we evaluate the cross-task transfer of LSC scaling factors across target test task $T$ and calibration task $C$.

As shown in Table~\ref{tab:cross_task}, LSC exhibits high task specificity. For instance, evaluating MNIST with CIFAR-10's parameters drops accuracy from \textbf{56.40\%} to \textbf{29.30\%}, proving that calibration profiles are task-specific. Interestingly, applying MNIST parameters to Fashion-MNIST test targets boosts accuracy to \textbf{55.70\%}, outperforming its native calibration (\textbf{50.00\%}). Since both share identical grayscale structures, they experience highly correlated variance collapse profiles (final block 4.1 scaling factors of 3.76 vs 3.83). MNIST's cleaner representation space yields a highly stable scaling sequence that generalizes exceptionally well to Fashion-MNIST, indicating transfer viability across topologically aligned domains.

\subsection{Computational and Storage Efficiency}\enlargethispage{3\baselineskip}
\label{sec:efficiency}
To validate the practical utility of our minimalist frameworks, we compare the calibration parameter storage (parameters per task), inference latency (ms per batch, size 128), and calibration time (ms to process $N=128$ samples on CPU) of all calibration methods. The benchmark results are summarized in Table~\ref{tab:efficiency}.

\begin{table}[h]
\caption{\textbf{Computational and Storage Efficiency} of calibration methods under WA ($N=128$). We report parameter storage (parameters per task), inference latency (ms per batch, size 128), and calibration time (ms for 128 samples).}
\label{tab:efficiency}
\vskip 0.05in
\begin{center}
\begin{small}
\begin{sc}
\setlength{\tabcolsep}{1.5pt}
\begin{tabular}{lccc}
\toprule
Method & Params/Task & Latency (ms) & Calib. (ms) \\
\midrule
NONE & \textbf{0} & \textbf{59.21} & --- \\
TCAC & 9,600 & 79.34 (+34.0\%) & 91.97 \\
SAC & 4,800 & 63.45 (+7.2\%) & 91.97 \\
LSC (Ours) & 20 & 63.47 (+7.2\%) & 91.97 \\
\textbf{TSC (Ours)} & \textbf{7} & \textbf{60.74 (+2.5\%)} & \textbf{91.97} \\
\bottomrule
\end{tabular}
\end{sc}
\end{small}
\end{center}
\vskip -0.1in
\end{table}

Our proposed LSC achieves a spectacular \textbf{480$\times$} reduction in calibration parameter storage compared to optimized TCAC (20 vs. 9,600 parameters), and our Threshold-gated TSC achieves a \textbf{1,371$\times$} parameter storage reduction (only 7 parameters per task) by calibrating only the most severely collapsed layers. Furthermore, while TCAC introduces a heavy \textbf{34.0\%} computational latency overhead during inference (79.34~ms vs. 59.21~ms), TSC is practically free, introducing only a negligible \textbf{2.5\%} latency overhead (60.74~ms). Calibration is extremely rapid, requiring only \textbf{91.97~ms} to process 128 samples, highlighting the high efficiency and practical value of our minimalist frameworks.

\section{Discussion \& Conclusion}\enlargethispage{2\baselineskip}
\label{sec:conclusion}
We demonstrated that complex, channel-wise affine methods (like TCAC) suffer from fatal numerical instabilities due to ReLU sparsity on small calibration batches. To resolve this, we proposed \textbf{Layer-wise Scaling-only Calibration (LSC)}, a stable, minimalist framework applying a single positive scalar per layer globally. LSC completely dominates more complex methods, where selectively calibrating early or late layers yields an outstanding \textbf{49.63\%} average accuracy (a \textbf{+8.3\% absolute boost}). By honoring the underlying sparse structure of neural networks and stripping away redundant components, we achieve robust and elegant model merging.

\textbf{Limitations \& Future Work}: We identify three limitations: (1) we focus on CNNs, and extending LSC to Transformers with non-ReLU activations is a vital future path; (2) LSC requires task identity during inference; and (3) downstream calibration cannot recover representation spaces that are completely destroyed at the parameter level (e.g., high drop rates in TIES/DARE).

\bibliography{paper}
\bibliographystyle{icml2026}

\clearpage
\appendix
\onecolumn
\section{Evaluating PLM-Centric Merging (TIES and DARE) on Convolutional Architectures}
\label{sec:appendix_advanced_merging}

While parameter-space merging techniques such as TIES-Merging~\cite{yadav2023ties} and DARE-Merging~\cite{yu2024dare} have achieved great success in Pre-trained Language Models (PLMs) using Transformer architectures, their performance and compatibility on Convolutional Neural Networks (CNNs) containing Batch Normalization (BatchNorm) layers remain largely unexplored. 

In this appendix, we conduct a systematic empirical evaluation of TIES and DARE on our ResNet-18 multi-task benchmark. We identify a fundamental limitation of these PLM-centric methods when applied to CNNs and demonstrate how our Threshold-gated Selective Calibration (TSC) interacts with them.

\subsection{Categorical Mismatch and Immediate Representation Collapse}
During our initial experiments, we observed that applying standard TIES or DARE directly to the model's parameters causes an immediate, catastrophic representation collapse to near-random guessing ($9.53\%$ average accuracy), even with extremely mild parameter modifications (e.g., dropping only $1\%$ of the task vector parameters in DARE).

We discover that this is driven by a subtle but severe category-level mismatch in the merging pipeline. Standard TIES and DARE implementations typically prune or mask all floating-point tensors in the model's state dictionary. In CNNs, this includes not only the weights and biases of convolutional layers but also the running statistics of Batch Normalization layers: \texttt{running\_mean} and \texttt{running\_var}. 

When a random mask is applied to these trackers, it corrupts the pre-calculated running statistics of BatchNorm, shifting and rescaling feature distributions arbitrarily and causing immediate network-wide signal collapse.

\subsection{Robust Sweeps with Tracker-Exclusion}
To evaluate the true capability of TIES and DARE, we implement a robust tracker-exclusion pipeline. Under this pipeline, BatchNorm running statistics and tracker buffers are excluded from pruning/dropping operations and are instead combined using simple, stable Weight Averaging (WA). Only the actual learnable weights and biases are subjected to TIES and DARE.

Using this robust pipeline, we sweep DARE drop rates ($d_r \in \{0.01, 0.05, 0.10, 0.20, 0.50\}$) and TIES trim fractions ($t_f \in \{0.99, 0.95, 0.90, 0.80, 0.50, 0.20\}$, representing the fraction of parameters retained). The multi-task accuracies under NONE (uncalibrated), TCAC, and our TSC ($\tau = 1.30$) are summarized in Table~\ref{tab:dare_sweep} and Table~\ref{tab:ties_sweep}.

\begin{table}[h]
\caption{\textbf{DARE-Merging Sweeps (\%) with BatchNorm Tracker-Exclusion} under varying drop rates $d_r$. We report task accuracies and multi-task averages for NONE, TCAC, and our TSC ($\tau = 1.30$).}
\label{tab:dare_sweep}
\vskip 0.1in
\begin{center}
\begin{small}
\begin{sc}
\begin{tabular}{lccccc}
\toprule
Method & Drop Rate $d_r$ & MNIST & Fashion-MNIST & CIFAR-10 & \textbf{Average} \\
\midrule
DARE + NONE & 0.01 & 49.00 & 31.10 & 27.40 & 35.83 \\
DARE + TCAC & 0.01 & 18.80 & 10.60 & 9.60 & 13.00 \\
\textbf{DARE + TSC (Ours)} & \textbf{0.01} & \textbf{67.50} & \textbf{50.50} & \textbf{27.40} & \textbf{48.47} \\
\midrule
DARE + NONE & 0.05 & 10.90 & 6.00 & 9.00 & 8.63 \\
DARE + TCAC & 0.05 & 9.30 & 10.20 & 9.40 & 9.63 \\
DARE + TSC (Ours) & 0.05 & 10.90 & 6.00 & 9.00 & 8.63 \\
\midrule
DARE + NONE & 0.10 & 7.60 & 10.10 & 11.30 & 9.67 \\
DARE + TCAC & 0.10 & 10.30 & 10.00 & 12.10 & 10.80 \\
DARE + TSC (Ours) & 0.10 & 7.60 & 10.10 & 11.30 & 9.67 \\
\midrule
DARE + NONE & 0.20 & 9.00 & 10.60 & 10.00 & 9.87 \\
DARE + TCAC & 0.20 & 11.00 & 10.90 & 9.30 & 10.40 \\
DARE + TSC (Ours) & 0.20 & 9.00 & 10.60 & 10.00 & 9.87 \\
\midrule
DARE + NONE & 0.50 & 10.20 & 8.10 & 9.50 & 9.27 \\
DARE + TCAC & 0.50 & 11.00 & 10.90 & 10.80 & 10.90 \\
DARE + TSC (Ours) & 0.50 & 10.20 & 8.10 & 9.50 & 9.27 \\
\bottomrule
\end{tabular}
\end{sc}
\end{small}
\end{center}
\end{table}

\begin{table}[h]
\caption{\textbf{TIES-Merging Sweeps (\%) with BatchNorm Tracker-Exclusion} under varying trim fractions $t_f$ (parameter retention fraction). We report task accuracies and multi-task averages for NONE, TCAC, and our TSC ($\tau = 1.30$).}
\label{tab:ties_sweep}
\vskip 0.1in
\begin{center}
\begin{small}
\begin{sc}
\begin{tabular}{lccccc}
\toprule
Method & Retained Fraction $t_f$ & MNIST & Fashion-MNIST & CIFAR-10 & \textbf{Average} \\
\midrule
TIES + NONE & 0.99 & 35.30 & 37.70 & 25.70 & 32.90 \\
TIES + TCAC & 0.99 & 19.30 & 11.80 & 10.00 & 13.70 \\
\textbf{TIES + TSC (Ours)} & \textbf{0.99} & \textbf{35.30} & \textbf{37.70} & \textbf{25.70} & \textbf{32.90} \\
\midrule
TIES + NONE & 0.95 & 10.60 & 11.40 & 17.10 & 13.03 \\
TIES + TCAC & 0.95 & 7.80 & 10.90 & 14.40 & 11.03 \\
TIES + TSC (Ours) & 0.95 & 10.60 & 11.40 & 17.10 & 13.03 \\
\midrule
TIES + NONE & 0.90 & 9.10 & 12.60 & 12.80 & 11.50 \\
TIES + TCAC & 0.90 & 11.10 & 9.00 & 10.80 & 10.30 \\
TIES + TSC (Ours) & 0.90 & 9.10 & 12.60 & 12.80 & 11.50 \\
\midrule
TIES + NONE & 0.80 & 9.10 & 10.00 & 9.00 & 9.37 \\
TIES + TCAC & 0.80 & 11.00 & 9.00 & 9.60 & 9.87 \\
TIES + TSC (Ours) & 0.80 & 9.10 & 10.00 & 9.00 & 9.37 \\
\midrule
TIES + NONE & 0.50 & 9.10 & 10.90 & 9.40 & 9.80 \\
TIES + TCAC & 0.50 & 11.00 & 9.00 & 9.60 & 9.87 \\
TIES + TSC (Ours) & 0.50 & 9.10 & 10.90 & 9.40 & 9.80 \\
\midrule
TIES + NONE & 0.20 & 9.10 & 9.70 & 9.60 & 9.47 \\
TIES + TCAC & 0.20 & 11.00 & 10.90 & 9.60 & 10.50 \\
TIES + TSC (Ours) & 0.20 & 9.10 & 9.70 & 9.60 & 9.47 \\
\bottomrule
\end{tabular}
\end{sc}
\end{small}
\end{center}
\end{table}

\subsection{Key Insights \& Scientific Findings}
Our robust sweeps reveal several critical, fundamental insights:

\textbf{1. Catastrophic Sensitivity to Weight Pruning}: Even with robust tracker-exclusion, CNNs remain incredibly fragile under weight pruning or dropping. Under DARE-Merging, dropping only $5\%$ of the weight parameters ($d_r = 0.05$) triggers immediate, irreversible representation collapse ($8.63\%$ average accuracy). Under TIES-Merging, pruning only $5\%$ of the weights ($t_f = 0.95$) degrades performance to $13.03\%$. At standard hyper-parameters (e.g., $d_r = 0.5$ or $t_f = 0.2$), both methods collapse completely to near-random guessing.

\textbf{2. Why CNNs are More Fragile than PLMs under Weight Pruning}:
\begin{itemize}
    \item \textbf{BatchNorm-Weight Mismatch}: In CNNs, BatchNorm running statistics are deeply coupled with the magnitudes and patterns of convolutional weights. Pruning weights, even slightly, changes the scale and distribution of feature maps. This violates the pre-computed BatchNorm trackers, leading to severe representation misalignment that compounds layer-by-layer.
    \item \textbf{Dense Spatial Connectivity}: In CNNs, convolutional kernels (typically 3x3) are highly compact and carry dense spatial filters. Setting even a single parameter to zero radically alters the spatial frequency response and the activation norms. In contrast, PLMs operate on massive, highly redundant feed-forward and attention projection matrices (typically 4096x4096), where the law of large numbers provides exceptional robustness to random drops.
\end{itemize}

\textbf{3. Role of Activation Calibration (TSC)}:
When the parameter-space drop rate is extremely low ($d_r = 0.01$ in DARE), our proposed TSC successfully rescues the representation space, boosting average accuracy from $35.83\%$ to \textbf{48.47\%} (+12.64\% absolute improvement). However, once drop rates or pruning fractions exceed $5\%$, the representation space is completely destroyed at the parameter level, representing an absolute barrier that downstream activation calibration cannot heal.

These findings suggest that PLM-centric pruning algorithms cannot be trivially transferred to convolutional architectures. For CNNs, simple and elegant Weight Averaging remains the most robust parameter-space baseline, which our proposed LSC and TSC can then elegantly and reliably calibrate to achieve state-of-the-art multi-task performance.

\end{document}
